\documentclass{article}
% ---------- Codificación e idioma ----------
\usepackage[utf8]{inputenc}
\usepackage[T1]{fontenc}
\usepackage{csquotes} 

% ---------- Matemáticas ----------
\usepackage{amsmath, amssymb} 

% ---------- Tipografía ----------
\usepackage[ttdefault=true]{AnonymousPro}
\usepackage{parskip}
\usepackage[protrusion=true, expansion=false]{microtype}

% ---------- Geometría ----------
\usepackage[margin=2.5cm]{geometry}

% ---------- Listas ----------
\usepackage{enumitem}
\setlist[itemize]{leftmargin=*, itemsep=2pt, topsep=3pt}
\setlist[enumerate]{leftmargin=*, itemsep=2pt, topsep=3pt}

% ---------- Cajas de color ----------
\usepackage{tcolorbox}
\usepackage{xcolor}

% ---------- Gráficos ----------
\usepackage{graphicx}
\usepackage{float}

% ---------- Bibliografía ----------
\usepackage[backend=biber, style=ieee, sorting=none]{biblatex}
\addbibresource{references.bib}

% ---------- Hipervínculos (cargar al final) ----------
\usepackage[hidelinks]{hyperref}

% COLORES INSTITUCIONALES
\definecolor{unamazul}{RGB}{0, 83, 155}
\definecolor{unagris}{RGB}{240, 240, 240}

% COMANDOS PERSONALIZADOS
\newcommand{\captura}[1]{
    \begin{center}
        \includegraphics[width=0.7\textwidth]{imgs/#1}
    \end{center}
}

\newcommand{\codigo}[1]{\texttt{#1}}

% PORTADA
\newcommand{\portadapractica}[5]{
\begin{titlepage}
    \centering
    \vspace*{1cm}
    {\scshape\LARGE Universidad Nacional Autónoma de México\par}
    \vspace{0.4cm}
    {\scshape\large Facultad de Ingeniería\par}
    \vspace{2.5cm}
    {\color{unamazul}\rule{\textwidth}{1.5pt}\par}
    \vspace{0.8cm}
    {\Huge\bfseries #2\par}
    \vspace{0.8cm}
    {\color{unamazul}\rule{\textwidth}{1.5pt}\par}
    \vspace{1.5cm}
    {\Large\itshape #1\par}
    \vspace{1.2cm}
    {\large\textbf{Profesor:} #3\par}
    \vspace{2.5cm}
    {\large\bfseries Integrantes\par}
    \vspace{0.6cm}
    {\large #4\par}
    \vfill
    {\large #5\par}
    \vspace{1cm}
\end{titlepage}
}

\begin{document}

\portadapractica
    {Programación Orientada a Objetos}
    {Ejemplo Flyweight: Generador de Árboles}
    {Edgar Tista García}
    {Robles Cancino Arturo Andres\\
    Rodríguez Cárcamo Ditter Emmanuel\\
    Santamaría Viveros César Yurem}
    {2 de Noviembre, 2026}

\section{Diseño del ejemplo: bosque con el patrón Flyweight}

En este ejemplo se modela un bosque que contiene distintos árboles.
Cada árbol tiene características propias, como su tipo, color y altura.
Sin embargo, es posible que muchos árboles compartan exactamente las
mismas características, por lo que no es necesario crear un objeto
independiente para cada uno.

Para resolver este problema se utiliza el patrón de diseño Flyweight,
que permite compartir objetos que contienen información común y
mantener por separado los datos que cambian dependiendo de cada árbol.

\section{Cómo funciona el código}

\subsection{El Flyweight: \texttt{Arbol}}

La clase \texttt{Arbol} representa el objeto que puede compartirse entre
varios árboles del bosque. Esta clase almacena las características
comunes entre los árboles, como el tipo, el color y la altura.

Estas características forman parte del estado intrínseco, debido a que
se mantienen iguales para todos los árboles que comparten esa instancia.

Además, la clase cuenta con el método \texttt{dibujar(int x, int y)},
que recibe las coordenadas del árbol como parámetros. De esta manera,
la posición no necesita almacenarse dentro del objeto compartido.

\subsection{La fábrica: \texttt{ArbolFactory}}

La clase \texttt{ArbolFactory} se encarga de administrar y reutilizar
los objetos de tipo \texttt{Arbol}.

Para conseguirlo, utiliza un \texttt{HashMap}, donde almacena las
instancias creadas de acuerdo con sus características: tipo, color y
altura.

Cuando se solicita un árbol, la fábrica primero verifica si ya existe
un objeto con esas características. Si existe, devuelve la instancia
almacenada; de lo contrario, crea un nuevo objeto y lo guarda para
poder reutilizarlo posteriormente.

Esto evita crear múltiples objetos con información idéntica y permite
que varios árboles compartan una misma instancia.

\subsection{El cliente: \texttt{Main}}

La clase principal representa al cliente que utiliza la fábrica para
obtener los árboles.

En lugar de crear directamente objetos de tipo \texttt{Arbol}, solicita
las instancias mediante \texttt{ArbolFactory}. Posteriormente, utiliza
el método \texttt{dibujar()} enviando las coordenadas correspondientes
a cada árbol.

De esta forma, varios árboles pueden compartir el mismo objeto,
pero conservar posiciones diferentes dentro del bosque.

\subsection{Ejemplo de ejecución}

En la salida se puede observar el tipo, color y altura
de cada árbol, así como su posición dentro del bosque y se puede ver que habiendo 
árboles con las mismas características, se reutiliza la misma instancia de \texttt{Arbol}.

La siguiente captura muestra un ejemplo de la ejecución del programa:

\captura{ejecucion.png}

\section{Estado intrínseco y extrínseco}

En este ejemplo se identifican dos tipos de estado:

\begin{itemize}
    \item \textbf{Estado intrínseco:} corresponde al tipo, color y
    altura del árbol. Esta información se almacena dentro de la clase
    \texttt{Arbol} y puede compartirse entre diferentes árboles.

    \item \textbf{Estado extrínseco:} corresponde a las coordenadas
    $(x,y)$ de cada árbol. Esta información depende de la ubicación
    particular de cada árbol y se proporciona al método
    \texttt{dibujar()} cuando se necesita.
\end{itemize}

La separación de estos estados permite reutilizar los datos que no
cambian y mantener de manera independiente la información que depende
de la ubicación de cada árbol dentro del bosque.

\section{¿Por qué es adecuado aplicar Flyweight aquí?}

Por ejemplo, si existen cien árboles de tipo pino, con el mismo color
y altura, no es necesario almacenar cien objetos con exactamente la
misma información. La fábrica puede crear una sola instancia de
\texttt{Arbol} y devolver esa misma referencia cada vez que se solicite
un árbol con dichas características.

Las principales razones por las que este código aplica el patrón son:

\begin{enumerate}
    \item \textbf{Reutilización de objetos:} la fábrica permite
    compartir instancias que tienen las mismas características.

    \item \textbf{Reducción del consumo de memoria:} se evita crear
    objetos duplicados que almacenan información idéntica.

    \item \textbf{Separación de estados:} las características comunes
    se almacenan en \texttt{Arbol}, mientras que las coordenadas se
    proporcionan externamente.

    \item \textbf{Uso de una fábrica:} la clase
    \texttt{ArbolFactory} centraliza la creación y recuperación de
    los objetos compartidos.

    \item \textbf{Independencia de la posición:} un mismo objeto
    puede utilizarse para representar árboles ubicados en diferentes
    coordenadas.
\end{enumerate}

\section*{Conclusión}

Se considera que este ejemplo implementa el patrón de diseño
Flyweight porque permite reducir la cantidad de objetos creados
mediante la reutilización de instancias que comparten las mismas
características.

La clase \texttt{Arbol} almacena el estado intrínseco, compuesto por
el tipo, color y altura, mientras que las coordenadas representan
el estado extrínseco. Por otra parte,
\texttt{ArbolFactory} administra las instancias para evitar crear
objetos repetidos.

De esta manera, el programa demuestra cómo el patrón Flyweight
permite optimizar el uso de memoria cuando se necesita representar
una gran cantidad de elementos que comparten información.

\end{document}